\newpage
\begin{center}
  \textbf{\large 2. ДИЗАЙН И МЕТОДОЛОГИЯ СИСТЕМЫ РЕКОНСТРУКЦИИ ВИЗУАЛЬНЫХ СТИМУЛОВ}
\end{center}
\refstepcounter{chapter}
\addcontentsline{toc}{chapter}{2. ДИЗАЙН И МЕТОДОЛОГИЯ СИСТЕМЫ РЕКОНСТРУКЦИИ ВИЗУАЛЬНЫХ СТИМУЛОВ}

\section{Общая постановка задачи}

Настоящая работа посвящена реконструкции бинарных визуальных стимулов размером $6 \times 6$ пикселей по многоканальным сигналам электроэнцефалографии (electroencephalography, EEG), зарегистрированным в фазе воспоминания ранее предъявленного изображения. В отличие от задач классификации, где модель выбирает одну метку класса, здесь требуется восстановить пространственную структуру изображения. Поэтому задача формулируется как многовыходная бинарная реконструкция.

Пусть $X_i \in \mathbb{R}^{C \times T}$ обозначает EEG-эпоху для $i$-го примера, где $C$ - число каналов, а $T$ - число временных отсчётов. Целевой стимул задаётся бинарной матрицей:

$$
Y_i \in \{0,1\}^{6 \times 6}.
$$

Для обучения моделей изображение дополнительно представляется в виде вектора из 36 элементов:

$$
\mathbf{y}_i \in \{0,1\}^{36}.
$$

Требуется построить отображение:

$$
f_\theta: \mathbb{R}^{C \times T} \rightarrow [0,1]^{36},
$$

где $\theta$ - параметры модели, а каждый выходной элемент интерпретируется как вероятность активности соответствующего пикселя. После получения вероятностей выполняется бинаризация, в результате которой формируется реконструированное изображение:

$$
\hat{\mathbf{y}}_i \in \{0,1\}^{36}.
$$

Таким образом, методология работы направлена не на распознавание типа стимула, а на восстановление локальной пиксельной структуры визуального образа, воспроизводимого испытуемым мысленно.

\section{Используемый набор данных}

В работе используется ранее собранный EEG-датасет визуального воображения, подготовленный в рамках предшествующего проекта по реконструкции визуальных стимулов. Датасет включает записи EEG и связанные с ними eye-tracking данные, полученные в контролируемом эксперименте с бинарными изображениями размером $6 \times 6$ пикселей.

Данные разделены по экспериментальным фазам. Первая часть относится к этапу предъявления и запоминания стимулов, а вторая - к этапу воспоминания и визуального воображения (visual imagery). Основной задачей данной магистерской работы является реконструкция изображения по EEG в фазе воспоминания, поэтому в основной обучающей и тестовой выборке используются данные второго типа. Данные этапа предъявления не смешиваются с ними, поскольку относятся к другой когнитивной фазе протокола. Они могут использоваться только как вспомогательный материал для контроля или последующего сравнения этапов запоминания и воспоминания.

Структура данных организована по испытуемым и экспериментальным пробам. Для каждой пробы хранятся EEG- и EOG-записи в формате \texttt{.fif}, а также файл разметки \texttt{labels.json}, задающий соответствие между записью и целевым стимулом. Формат разметки включает список блоков, где для каждого блока указаны индекс execution-блока, тип стимула, идентификатор паттерна и бинарная матрица изображения размером $6 \times 6$. Более подробное описание файловой организации датасета и структуры разметки приведено в приложении А.

Согласно описанию датасета, записи получены у 16 испытуемых. Регистрация EEG выполнялась с использованием системы actiCHamp с 64 каналами, частотой дискретизации 1000 Hz и расположением электродов по схеме 10-20. Для контроля взгляда и последующего удаления окулярных артефактов использовался EyeLink 1000 Plus. Итоговый корпус включает 31 trial и 651 синхронизированный EEG/eye-tracking sample, из которых 434 относятся к геометрическим паттернам, а 217 - к случайным паттернам. Длина одного sample составляет 16 секунд: 15 секунд целевой эпохи воспоминания и дополнительное окно около $\pm 0.5$ секунды для учёта динамики до и после целевого интервала.

\section{Экспериментальный протокол}

Исходный экспериментальный протокол строился вокруг двух последовательных когнитивных этапов: предъявления стимула для запоминания и последующего мысленного воспроизведения этого стимула по команде. В фазе воспоминания участник сохранял открытые глаза и фиксировал взгляд на специальной точке, что позволяло контролировать глазодвигательные артефакты и снижать вариативность зрительного поведения.

В рамках одной сессии выполнялись три task block. Первые два блока включали геометрические паттерны, а третий блок - случайные паттерны. В каждом task block участник сначала запоминал три паттерна и связывал их с визуальными командами, после чего выполнял последовательность execution sub-block. В каждом execution sub-block команда предъявлялась в течение 3 секунд, затем следовала фаза мысленного воспроизведения длительностью 15 секунд, после чего выполнялся 10-секундный период отдыха. Полный протокол предусматривал 21 целевую эпоху на сессию и до 42 целевых эпох на испытуемого.

Стимулы представляли собой бинарные изображения размером $6 \times 6$ пикселей. Использование такого малого разрешения обусловлено ограниченным пространственным разрешением EEG и необходимостью сохранить задачу реконструкции экспериментально контролируемой. Изображения предъявлялись на нейтральном сером фоне, а фиксационная точка использовалась для стабилизации взгляда во время выполнения задания. В текущей работе эти особенности протокола учитываются при интерпретации данных, однако основная обработка выполняется уже на подготовленных EEG/EOG sample, соответствующих фазе воспоминания.

\section{Предобработка сигналов}

В данной работе используются подготовленные EEG-записи в формате \texttt{.fif}. Методологически preprocessing pipeline включает несколько этапов, направленных на повышение качества сигнала и согласование EEG с временной структурой эксперимента.

Сначала выполняется восстановление пропущенных участков, если они возникли из-за аппаратных задержек или потерь samples при записи. Затем raw EEG переводится в микровольты, после чего применяется average reference, то есть вычитание среднего потенциала по каналам. Это снижает вклад общего шума и делает пространственное распределение потенциалов более сопоставимым между каналами.

Далее применяется полосовая FIR-фильтрация в диапазоне 1-40 Hz. Нижняя граница удаляет медленные дрейфы, а верхняя граница снижает вклад высокочастотного шума и мышечных артефактов. Для удаления окулярных артефактов используется независимый компонентный анализ (Independent Component Analysis, ICA). После разложения сигнала на компоненты оценивается связь компонент с масками морганий и саккад, полученными из eye-tracking/EOG данных. Компоненты, связанные с глазодвигательной активностью, исключаются, после чего EEG-сигнал реконструируется из оставшихся компонент.

На завершающем этапе из непрерывной записи выделяются целевые эпохи, соответствующие фазе воспоминания. Именно эти эпохи используются для построения признаков, обучения моделей и оценки качества реконструкции.

\section{Признаковое представление}

Для построения моделей реконструкции рассматриваются несколько представлений EEG-эпохи. Общий вид отображения признаков задаётся как:

$$
\phi: \mathbb{R}^{C \times T} \rightarrow \mathbb{R}^{d},
$$

где $\phi$ - процедура извлечения признаков, а $d$ - размерность полученного признакового пространства.

Для классических baseline-моделей используются компактные признаки, связанные с пространственной и спектральной структурой EEG. В первую очередь рассматриваются ковариационные матрицы каналов, поскольку они отражают межканальные зависимости и согласуются с многоканальной природой EEG. Дополнительно могут использоваться признаки спектральной мощности в альфа- и бета-диапазонах, а также Morlet power как расширенное временно-частотное представление.

Для EEGNet и EEGNet-подобной модели входом может служить предобработанная EEG-эпоха или её оконное представление. В этом случае извлечение признаков частично переносится внутрь нейросетевой архитектуры: временные и пространственные фильтры обучаются совместно с задачей предсказания пикселей.

Целевое изображение во всех вариантах представляется вектором из 36 бинарных значений. Это позволяет использовать единый формат выхода как для независимых пиксельных классификаторов, так и для многовыходной нейросетевой модели.

\section{Модели реконструкции}

С учётом ограниченного объёма данных и срока выполнения работы стратегия моделирования строится от простых baseline-моделей к более сложным архитектурам.

Первую группу составляют простые baseline-подходы. Случайный baseline предсказывает значения пикселей независимо, majority baseline выбирает наиболее частое значение, а pixel-frequency baseline учитывает эмпирическую частоту активности каждого пикселя в обучающей выборке. Эти модели не используют EEG-сигнал напрямую, но задают нижнюю границу качества и позволяют проверить, превосходит ли реконструкция статистику распределения стимулов.

Вторую группу составляют классические модели машинного обучения. Основной безопасный вариант - независимые пиксельные классификаторы, где для каждого из 36 пикселей обучается отдельная Logistic Regression или Linear SVM. Такой подход прост, интерпретируем и устойчивее к малому объёму данных, чем тяжёлые генеративные модели. Он также позволяет анализировать, какие пиксели реконструируются устойчивее.

Основной нейросетевой моделью рассматривается EEGNet или компактная EEGNet-подобная сверточная архитектура. Такая модель использует EEG-эпоху как вход и формирует 36 выходных вероятностей. EEGNet выбран как реалистичный компромисс между выразительностью и контролируемой сложностью: модель способна обучать пространственно-временные фильтры, но остаётся существенно легче, чем крупные transformer-based или generative модели.

BrainBERT рассматривается как дополнительный современный эксперимент. Если готовая реализация и предобученные веса могут быть быстро адаптированы к текущему формату данных, модель может использоваться как feature extractor или дообучаемый encoder. При этом BrainBERT не является единственной центральной моделью работы, поскольку его успешная адаптация зависит от совместимости данных, вычислительных ресурсов и риска переобучения.

Генеративные подходы, такие как GAN и diffusion models, не рассматриваются как основное направление текущей работы. Для имеющегося объёма данных они несут высокий риск нестабильного обучения, усреднённых реконструкций и слабой интерпретируемости. Их роль ограничивается обсуждением в обзоре литературы и возможными направлениями дальнейших исследований.

\section{Стратегии обучения и валидации}

Основной протокол обучения строится на данных фазы воспоминания. Данные фазы предъявления и запоминания не добавляются в основную обучающую выборку, поскольку это привело бы к смешению разных когнитивных состояний. Такое разделение важно для сохранения содержательной постановки задачи: модель должна восстанавливать изображение по EEG, связанному именно с visual imagery.

В работе рассматриваются несколько стратегий оценки. Внутри-субъектная оценка (intra-subject или within-subject) предполагает обучение и тестирование на данных одного испытуемого при корректном разделении эпох. Этот сценарий показывает, насколько модель может адаптироваться к индивидуальным пространственно-спектральным особенностям EEG.

Межсубъектная оценка (inter-subject или cross-subject) предполагает обучение на группе испытуемых и тестирование на данных другого испытуемого или другой группы. Такой сценарий сложнее, но он важен для проверки переносимости модели. При достаточном объёме экспериментов дополнительно может использоваться subject-independent схема, например leave-one-subject-out, где один испытуемый полностью исключается из обучения и используется только для тестирования.

При любом варианте валидации необходимо контролировать утечки данных. Примеры, относящиеся к одной и той же пробе или одному и тому же trial, не должны попадать в конфликтующие части выборки, если выбранный протокол предполагает независимость train и test. Это особенно важно при малом числе trials, когда случайное разбиение может завысить качество модели.

\section{Метрики оценки}

Качество реконструкции оценивается на уровне пикселей и на уровне формы восстановленного бинарного изображения. Для пиксельной постановки используются стандартные классификационные метрики precision, recall и F1-score~\cite{10.1016/j.ipm.2009.03.002}, overlap-метрики для сравнения бинарных масок~\cite{10.1186/s12880-015-0068-x}, а также Hamming loss, связанный с расстоянием Хэмминга~\cite{10.1002/j.1538-7305.1950.tb00463.x}. Формальные определения метрик приведены в приложении~\ref{app:metrics}. Основными метриками являются:

\begin{itemize}
\item Pixel Accuracy - доля правильно восстановленных пикселей;
\item F1-score - баланс между precision и recall для активных пикселей;
\item Intersection over Union (IoU) - степень совпадения активных областей истинного и реконструированного изображения;
\item Hamming loss - доля пикселей, предсказанных неверно.
\end{itemize}

Дополнительно может использоваться balanced accuracy по пикселям, если распределение активных и неактивных пикселей существенно несбалансировано. Exact match accuracy, то есть доля полностью восстановленных изображений, рассматривается только как строгая вспомогательная метрика, поскольку для изображения из 36 пикселей она может быть чрезмерно жёсткой.

Все метрики должны интерпретироваться относительно baseline-моделей. Сравнение с random, majority и pixel-frequency baseline необходимо для того, чтобы отделить реальное использование EEG-информации от воспроизведения статистики обучающего набора.

\section{Ограничения методологии}

Методология работы учитывает несколько ограничений. Во-первых, EEG обладает низким пространственным разрешением, поэтому восстановление пиксельной структуры изображения является сложной и шумной задачей. Во-вторых, визуальное воображение менее стабильно, чем прямое зрительное восприятие: испытуемые могут использовать разные внутренние стратегии воспроизведения образа.

В-третьих, объём данных ограничен для обучения тяжёлых моделей. Это повышает риск переобучения, особенно при использовании transformer-based или generative architectures. По этой причине центральное место в экспериментальной стратегии занимают классические baseline-модели и EEGNet, а BrainBERT рассматривается как дополнительный эксперимент при наличии технической возможности.

В-четвёртых, возможен дисбаланс активных и неактивных пикселей. Если отдельные области изображения встречаются чаще других, модель может склоняться к предсказанию усреднённых или наиболее частых паттернов. Поэтому оценка качества должна обязательно включать baseline-сравнение и метрики, чувствительные к активным пикселям.

Несмотря на эти ограничения, выбранная методология позволяет сформулировать воспроизводимую задачу реконструкции визуального образа малого разрешения по EEG в фазе воспоминания и сравнить несколько моделей разной сложности в едином экспериментальном протоколе.
